This dissertation presents, implements, and validates an automated and scalable methodology for assessing performance on Sustainable Development Goals (SDGs) through natural language processing (NLP) and large language models (LLMs). The methodology combines three main components: automated collection of public reviews via the Google Places API, hierarchical semantic filtering with multiple geographic and content validation criteria, and aspect-based sentiment analysis using Gemini and Perplexity models. The system was applied to the case study of Mossoró-RN, the second largest city in the state of Rio Grande do Norte (Brazil), which presents very low scores (0--39.99) on SDG 9 (Industry, Innovation and Infrastructure) and SDG 16 (Peace, Justice and Strong Institutions). A total of 253 reviews were collected and analyzed for SDG 9 (117 establishments) and 116 for SDG 16 (12 establishments), with average confidence scores above 0.84. Results revealed that 56.9\% of SDG 9 analyses identify problems, concentrated in physical infrastructure degradation (especially the bus terminal), with cascading impacts on mobility and local entrepreneurship. For SDG 16, 46.6\% of analyses report problems, dominated by public institutional inefficiency (37.1\% of mentions), which amplifies barriers to access to justice and public safety. The system also generates structured natural language diagnostics and feeds interactive Power BI dashboards, making results accessible to public managers without advanced technical expertise. The modular architecture based on JSON configuration files enables immediate extension to any SDG and municipality, and is made publicly available as a contribution to the scientific and public management communities.

% Separe as Keywords por ponto
\keywords{Sustainable Development Goals. Natural Language Processing. Sentiment Analysis. Large Language Models. Google Places API. Public Policy. Mossoró-RN.}
